\section{Q5: Graph Measures (30 pts)}

All measures below are computed on the graph built in Q3: the simple undirected
graph read from \texttt{graph.txt} (one \texttt{networkx.Graph}, one node per
distinct id, one edge per line), which has $26{,}728$ nodes, $26{,}377$ edges and
$4104$ connected components, the largest of which contains $13{,}903$ nodes and
$16{,}695$ edges. All computations use NetworkX with its default parameters; every
number was reproduced under NetworkX 3.3 and 3.6.1 and independently re-derived
with SciPy (see the remarks).

\subsection{Eigenvector centrality (10 pts)}

\texttt{nx.eigenvector\_centrality(G)} was called with default parameters
(\texttt{max\_iter}$=100$, \texttt{tol}$=10^{-6}$). \textbf{It converged}; no
\texttt{PowerIterationFailedConvergence} was raised, so the default call is the
one reported. Values are the entries of the principal eigenvector of the adjacency
matrix, normalised to unit $L^2$ norm.

\begin{table}[H]
\centering
\begin{tabular}{clrcc}
\toprule
Rank & Node & Degree & Centrality (default) & Centrality (exact) \\
\midrule
1  & 16485 & 57 & 0.491844 & 0.492371 \\
2  & 22777 & 59 & 0.419863 & 0.422581 \\
3  & 14618 & 38 & 0.177656 & 0.177494 \\
4  & 21928 & 17 & 0.139183 & 0.139695 \\
5  & 2884  & 29 & 0.133813 & 0.133861 \\
6  & 1197  & 18 & 0.133186 & 0.133738 \\
7  & 9068  & 51 & 0.111138 & 0.102654 \\
8  & 17902 & 32 & 0.109219 & 0.109103 \\
9  & 1324  & 20 & 0.100115 & 0.100053 \\
10 & 25747 & 19 & 0.094368 & 0.094399 \\
\bottomrule
\end{tabular}
\label{p3:tab:eig}
\end{table}

\paragraph{Remark on the default tolerance.}
The fourth column is the answer with default parameters, as asked. For comparison
the fifth column is the exact eigenvector from
\texttt{eigenvector\_\-centrality\_\-numpy}. The two agree on the \emph{set} of
top-10 nodes, but they order ranks 7 and 8 differently: the exact computation puts
node 17902 ahead of node 9068. The reason is that NetworkX's power iteration stops
when the total $L^1$ change falls below
\[
  n\cdot\texttt{tol} \;=\; 26{,}728\times10^{-6} \;\approx\; 2.67\times10^{-2},
\]
which is a loose \emph{absolute} bound on a vector of unit norm; the returned
vector has residual $\lVert Ax-\lambda x\rVert_\infty\approx7.8\times10^{-3}$ and
Rayleigh quotient $8.76202$ against the true
$\lambda_{\max}=8.762217$. Node 9068 is the entry still
furthest from its limit, so its rank has not settled after 100 iterations. Running
the same power iteration with \texttt{max\_iter}$=10^5$, \texttt{tol}$=10^{-10}$
reproduces the exact column to $9.3\times10^{-7}$.

\paragraph{Interpretation.}
Because the graph is disconnected, the principal eigenvector is supported on the
component with the largest spectral radius. Here that is the giant component
($\lambda_{\max}=8.7622$, versus at most $3.63$ for every other component), so all
ten top nodes lie in the giant component and roughly half the nodes
($12{,}908$, or $48.3\%$) receive centrality below $10^{-12}$. Note also that
eigenvector centrality is not simply degree: only 6 of these 10 nodes are also in
the top 10 by degree. Node 21928 ranks 4th with degree just 17, because its
neighbours are themselves central, while node 17546 has degree 40 yet does not
make the list.

\subsection{Katz centrality (10 pts)}

\texttt{nx.katz\_centrality(G)} was called with default parameters
($\alpha=0.1$, $\beta=1.0$, \texttt{max\_iter}$=1000$, \texttt{tol}$=10^{-6}$),
solving $x=\beta\mathbf{1}+\alpha Ax$ and normalising to unit $L^2$ norm.

\paragraph{Admissibility of the default $\alpha$.}
Katz centrality is well defined only for $\alpha<1/\lambda_{\max}$. For this graph
$\lambda_{\max}=8.762217$, so $1/\lambda_{\max}=0.114126$. The default
$\alpha=0.1$ satisfies $0.1<0.114126$, so it \textbf{is} admissible (though not by
a large margin). Consistently with this, the default call \textbf{converged} and
every entry of the unnormalised solution is positive.

\begin{table}[H]
\centering
\begin{tabular}{clrc}
\toprule
Rank & Node & Degree & Katz centrality \\
\midrule
1  & 16485 & 57 & 0.211609 \\
2  & 22777 & 59 & 0.175512 \\
3  & 9068  & 51 & 0.109482 \\
4  & 14618 & 38 & 0.099380 \\
5  & 2884  & 29 & 0.074102 \\
6  & 17902 & 32 & 0.070590 \\
7  & 17546 & 40 & 0.068771 \\
8  & 21928 & 17 & 0.066834 \\
9  & 1197  & 18 & 0.061642 \\
10 & 9740  & 27 & 0.060545 \\
\bottomrule
\end{tabular}
\label{p3:tab:katz}
\end{table}

These values were confirmed against \texttt{nx.katz\_centrality\_numpy(G)} and
against a direct sparse solve of $(I-\alpha A)x=\beta\mathbf{1}$; all three give
the same ordering, agreeing to $1.3\times10^{-5}$ and $1.9\times10^{-16}$
respectively. Unlike eigenvector centrality, the $\beta\mathbf{1}$ term gives every
node a positive baseline, so Katz does not collapse to zero off the giant
component, and its ranking tracks degree more closely (node 9068, degree 51, rises
from 7th to 3rd).

\subsection{Average clustering coefficient and average Jaccard similarity (10 pts)}

\paragraph{Average clustering coefficient.}
Using \texttt{nx.average\_clustering(G)}, i.e.\ the mean over \emph{all} $26{,}728$
nodes of the local clustering coefficient
$c_u = 2T_u / \bigl(d_u(d_u-1)\bigr)$, where $T_u$ is the number of triangles
through $u$, with the standard convention that $c_u=0$ for nodes of degree $<2$:
\[
  \bar{C} \;=\; 0.100539 .
\]
This average is held down by the graph's sparsity: $13{,}942$ of the $26{,}728$
nodes ($52.2\%$) have degree $<2$ and so contribute $0$ by convention, and the
graph contains only $1926$ triangles in total. Restricted to the $12{,}786$ nodes
of degree $\ge 2$ the mean local clustering is $0.210168$.

\paragraph{Average Jaccard similarity of all connected node pairs.}
For every edge $(u,v)\in E$ the Jaccard similarity of the two neighbourhoods is
\[
  J(u,v) \;=\; \frac{\lvert N(u)\cap N(v)\rvert}{\lvert N(u)\cup N(v)\rvert},
\]
computed with \texttt{nx.jaccard\_coefficient(G, G.edges())} and averaged over all
$26{,}377$ edges. Here $N(\cdot)$ denotes the \emph{raw} neighbour set, so
$v\in N(u)$ and $u\in N(v)$; consequently an isolated edge (both endpoints of
degree 1) has $\lvert N(u)\cap N(v)\rvert=0$ and scores $J=0$ rather than being
undefined. Under this convention
\[
  \bar{J} \;=\; 0.041304 .
\]
The value is small because $21{,}167$ of the $26{,}377$ edges ($80.3\%$) join two
nodes with no common neighbour and therefore score exactly $0$; the largest value
attained by any edge is $2/3$. (For contrast, excluding the two endpoints from each
other's neighbour sets would give $0.084305$; the number reported above is the
NetworkX convention stated in the formula.)
